% !PRINTED 640
\weberpage{640}
%<<CONT

und erhalten

\weq{5}{(\varepsilon^{r_\nu}, \xi) = K_\nu \tau_\nu^{r_{n-2}} \tau_{\nu+1}^{r_{n-3}} \ldots \tau_{\nu+n-2}^{r_0},}

und wenn man diese Formeln alle addirt, und noch die rationale Grösse $(1, x) = A$ setzt:

\weq{6}{n\xi_0 = A + \sum_{0,\,n-2}^{\nu} K_\nu \tau_\nu^{r_{n-2}} \tau_{\nu+1}^{r_{n-3}} \ldots \tau_{\nu+n-2}^{r_0}.}

Nach \S.\,184, (16) und (17) sind die Grössen $K_\nu$ und $k_\nu$ alle von Null verschieden.

Durch (6) ist eine der Wurzeln $\xi$ dargestellt.

Setzen wir zur Abkürzung

\[
R_\nu = k_\nu^{r_{n-2}} k_{\nu+1}^{r_{n-3}} \ldots k_{\nu+n-2}
\]

und schreiben (6) in der Form

\weq{7}{n\xi_0 = A + K_0 \sqrt[n]{R_0} + K_1 \sqrt[n]{R_1} + \cdots + K_{n-2} \sqrt[n]{R_{n-2}},}

so ist also $\xi_0$ durch $n-1$ Radicale $n^{\mathrm{ten}}$ Grades ausgedrückt, von denen jedes an sich $n$ verschiedene Werthe haben kann. Diese Werthe können aber nicht von einander unabhängig sein, weil sonst die Anzahl der Werthe von $\xi$, die sich aus (7) ergeben, zu gross wäre. Man erhält in der That aus (5), wenn man nach \S.\,184, (16)

\[
(\varepsilon^{r_\nu}, \xi) (\varepsilon^{r_{\nu-1}}, \xi)^{-g} = k_{\nu-1}
\]

setzt,

\weq{8}{K_\nu \sqrt[n]{R_\nu} = K_{\nu-1}^g k_{\nu-1} \left( \sqrt[n]{R_{\nu-1}} \right)^g,}

und kann also hiernach alle diese Radicale rational durch eines von ihnen und durch $k_\nu$ ausdrücken.

Der Ausdruck (6) hat aber vor (7) den grossen Vorzug, \gesperrt{dass er, wie man auch die $n-1$ Radicale $\sqrt[n]{k_\nu}$ bestimmen mag, doch nur $n$ verschiedene Werthe, nämlich die $n$ Wurzeln $\xi$ darstellt}.

Um dies nachzuweisen, bezeichnen wir mit $\varepsilon_0, \varepsilon_1 \ldots \varepsilon_{n-2}$ irgend ein beliebiges System $n^{\mathrm{ter}}$ Einheitswurzeln, und ersetzen in (5)

\[
\sqrt[n]{k_0}, \quad \sqrt[n]{k_1} \ldots \sqrt[n]{k_{n-2}}
\]

durch

\[
\varepsilon_0 \sqrt[n]{k_0}, \; \varepsilon_1 \sqrt[n]{k_1} \ldots \varepsilon_{n-2} \sqrt[n]{k_{n-2}}:
\]
%CONT>>
